% ============================================================================ CONCLUSION
\section{Conclusion}\label{sec:conclusion}

Integrating a voltage regulator at the point of load is not a placement choice that can be made
after the converter is designed: the integration sets the converter's temperature and therefore
its loss, the waveform its package interconnect carries, the deformation of the package and the
survival of the interconnect and ball-field rules. The framework presented here carries a
synthesized hybrid switched-capacitor--buck regulator through an explicit package to these
consequences on one geometry, verifies each stage against a matched reference, and does so in
\EvalHours{} per integration on a shared workstation. Used to evaluate, compare and explore, it showed on one 1-kW
specification that coupling reverses the electrical ranking of the integrations in both
directions, that the embedded regulator requires local input decoupling and a repartitioned tier
order, that the substrate-backside regulator warps the die \WarpRatioZtwo{} times as much as the
embedded one, that the board-backside regulator is limited by the ball field rather than by its
thermal or mechanical behavior, and that the regulator-count optimum shifts with the topology
chosen at each count.
The same chain applies to any member of the hybrid family, any package stack the builder can
express, and any cooling solution with a declared boundary, which makes it a design tool for the
integration decision itself rather than a check applied after it.
